\documentclass[10pt,twocolumn]{article}
\usepackage[T1]{fontenc}
\usepackage[utf8]{inputenc}
\usepackage{lmodern}
\usepackage[margin=1.8cm,top=2.4cm,bottom=2cm,columnsep=0.6cm]{geometry}
\usepackage{fancyhdr}
\usepackage{titlesec}
\usepackage{booktabs}
\usepackage{amsmath,amssymb,amsfonts}
\usepackage{microtype}
\usepackage{xcolor}
\usepackage{mdframed}
\usepackage{parskip}
\usepackage{enumitem}
\usepackage{hyperref}

\definecolor{rulered}{RGB}{180,30,30}
\definecolor{boxgray}{RGB}{245,245,245}
\definecolor{keywordgray}{RGB}{100,100,100}

\pagestyle{fancy}
\fancyhf{}
\fancyhead[L]{\textsc{\small Claude Mag}}
\fancyhead[C]{\small\textit{Machine Learning Research Digest}}
\fancyhead[R]{\small 2026-09-30}
\fancyfoot[C]{\small\thepage}
\renewcommand{\headrulewidth}{0.8pt}

\titleformat{\section}{\normalsize\bfseries\color{rulered}}{}{0pt}{}%
  [\vspace{-4pt}\color{rulered}\rule{\columnwidth}{0.4pt}]
\titlespacing{\section}{0pt}{8pt}{4pt}

\hypersetup{colorlinks=true,urlcolor=rulered,linkcolor=black}

\begin{document}
\setlength{\parindent}{0pt}
\setlength{\parskip}{4pt}

{\small\color{keywordgray}\textbf{Keywords:}
  KV cache \textbullet{}
  sparse attention \textbullet{}
  long-context inference \textbullet{}
  attention compensation}\par\vspace{4pt}

{\LARGE\bfseries\raggedright
  CompKV}\par\vspace{4pt}

{\large\itshape\raggedright
  Compensation-Aware KV Selection for Long-Context LLM Inference}\par\vspace{6pt}

{\small\textbf{By} Anonymous et al.}\par\vspace{2pt}

{\small\color{keywordgray}arXiv:2609.26300 \textbullet{} September 2026}
\par\vspace{2pt}\noindent{\color{rulered}\rule{\columnwidth}{0.6pt}}\vspace{6pt}

By pairing block-level KV selection with mean compensation---which yields an exact
log-partition identity---CompKV derives a selection priority proportional to block
attention mass times within-block logit variation, consistently outperforming
prior sparse-attention methods on RULER and LongBench-Pro at tight token budgets.

\begin{mdframed}[backgroundcolor=boxgray,linecolor=rulered,linewidth=1pt,
    innerleftmargin=6pt,innerrightmargin=6pt,
    innertopmargin=6pt,innerbottommargin=6pt]
\textbf{\small AT A GLANCE}\par\vspace{4pt}
\begin{description}[leftmargin=0pt,labelsep=4pt,itemsep=2pt,parsep=0pt,
    font=\small\bfseries]
  \item[Problem:] Sparse-attention KV selection methods discard unselected
    context entirely, introducing approximation error that biases the softmax
    partition function.
  \item[Why it matters:] Long-context tasks requiring global aggregation fail
    most under naive selection; more principled compensation preserves performance
    at aggressive memory budgets.
  \item[Method:] Mean compensation for unselected blocks yields an exact
    log-partition identity; the resulting optimal selection priority is
    block attention mass $\times$ within-block logit variation.
  \item[Result:] CompKV consistently beats selection-only and eviction-based
    sparse attention on Llama-3.1-8B and Qwen3 models.
  \item[Evidence:] RULER and LongBench-Pro benchmarks at context lengths
    up to 128K tokens.
\end{description}
\end{mdframed}

\section{The Big Picture}

Long-context LLM inference is memory-bound. Fitting a 128K-token context in a
transformer requires gigabytes of key-value (KV) cache that exceed typical GPU
memory. Sparse-attention methods reduce pressure by loading only a selected
subset of KV blocks per query. The dominant selection heuristic---loading the
top-$K$ blocks by attention mass---has a hidden flaw: it treats unselected
blocks as if they carry zero weight, silently biasing the softmax partition
function and the attention output.

CompKV revisits this approximation and asks: given that some blocks will not
be loaded, what is the best way to account for them? The answer is mean
compensation---replacing each unselected block with its mean---combined with a
principled selection criterion derived from the compensation's residual error.

\section{Why Existing Approaches Fall Short}

\textbf{Eviction methods} (H2O, StreamingLLM) permanently discard tokens,
making them unrecoverable for later attention layers.

\textbf{Selection without compensation} loads the top-$K$ blocks and zeroes
the rest, which underestimates the softmax denominator whenever unselected
blocks carry non-trivial mass.

\textbf{KV quantization} reduces per-token size but does not reduce the number
of tokens read per attention call.

\textbf{IO-aware kernels} overlap memory transfers with compute but leave the
approximation error unaddressed.

None of these methods reason about how unselected context should be summarised
to correct the partition function.

\section{Core Mechanism}

\textbf{Mean Compensation.} For each unselected block $j$, CompKV replaces
every logit with the block mean $\bar{e}_j$. This is the unique single-number
summary that preserves the block's contribution to the log-partition function
to first order---a property that zero-padding and other approximations lack.

\textbf{Compensation-Aware Priority.} After fixing mean compensation, the
residual approximation error for block $j$ is bounded by the product of its
attention mass $m_j$ (how much it contributes to the output) and its within-block
logit standard deviation $\sigma_j$ (how much individual weights differ from the
mean). CompKV selects the blocks with the largest $m_j \cdot \sigma_j$---those
for which mean compensation is most inadequate.

\section{Technical Formulation}

Let the KV cache be partitioned into blocks of size $b$. For block $j$, denote
the attention logits $\mathbf{e}_j \in \mathbb{R}^b$ and the block mean
$\bar{e}_j = \frac{1}{b}\sum_i e_{j,i}$.

The log-partition under mean compensation:
\[
\log Z_j^{\text{approx}} = \bar{e}_j + \log b
\]
which equals the exact value when all within-block logits are equal. The
residual error is bounded:
\[
\Delta_j \;\leq\; m_j \cdot \sigma_j
\]
where $m_j = \mathrm{softmax\_mass}(\mathbf{e}_j)$ and
$\sigma_j = \mathrm{std}(\mathbf{e}_j)$. The CompKV selection priority score is:
\[
s_j = m_j \cdot \sigma_j
\]
The top-$K$ blocks by $s_j$ are loaded; the rest contribute via their mean.

\section{Learning or Inference Procedure}

CompKV is training-free and applies at inference time:

\begin{enumerate}[itemsep=2pt,parsep=0pt]
  \item A fused GPU kernel computes per-block mass estimates and logit
    standard deviations using grouped variance projection of the query
    against pre-computed block mean keys.
  \item Blocks are ranked by $s_j = m_j \cdot \sigma_j$; top-$K$ are selected.
  \item Selected blocks contribute via exact attention; unselected blocks
    contribute via their mean, with the partition function adjusted accordingly.
\end{enumerate}

\textbf{Implementation:} BF16 KV cache and FP32 block summaries are kept in
pinned CPU memory. A main CUDA stream handles selection-region transfer and
exact-attention kernels; an auxiliary stream transfers mean values in parallel,
eliminating blocking synchronization.

\section{What the Guarantee Says}

Mean compensation guarantees exact log-partition preservation (to first order)
when within-block logits are uniform. For non-uniform blocks, the residual
error is bounded by $m_j \cdot \sigma_j$---precisely the quantity used for
selection. Under a fixed budget of $K$ blocks, CompKV minimises the total
approximation error bound by choosing the blocks for which the bound is largest.

\section{Experimental Findings}

\begin{itemize}[itemsep=2pt,parsep=0pt]
  \item Models: Llama-3.1-8B-Instruct, Qwen3-8B, Qwen3-32B
  \item Benchmarks: RULER (controlled retrieval and aggregation), LongBench-Pro
    (realistic long-context tasks), context lengths to 128K tokens
  \item CompKV outperforms selection-only (top-$K$ by mass), H2O, and
    StreamingLLM variants consistently across all models and benchmarks
  \item Gains are largest on aggregation tasks and at tighter budgets ($K$
    is small), where the compensation correction matters most
\end{itemize}

\section{Ablations and Interpretation}

\textbf{Selection criterion:} Replacing $m_j \cdot \sigma_j$ with mass-only
or variation-only degrades performance, confirming that the joint criterion is
necessary.

\textbf{Compensation type:} Mean compensation outperforms zero compensation;
the gap grows with block size (larger within-block variation).

\textbf{IO overhead:} The fused kernel architecture hides transfer latency by
overlapping streams; throughput is competitive with GPU-only sparse attention
at high batch sizes.

\section*{Reference}

Anonymous et al.
\textbf{CompKV: Compensation-Aware KV Selection for Long-Context LLM Inference}.
arXiv:2609.26300, September 2026.\\
\url{https://arxiv.org/abs/2609.26300}

\end{document}
